% Transcription — fonds Alexandre Grothendieck, shelfmark 7, batch 2 (pages 21-40).
% Established from the facsimile archives/batches/7/batch-02.pdf, cut from a
% copy of 7.pdf fetched through the project's own relay
% (https://grothendieck-relay.onrender.com/source/7.pdf) rather than from
% grothendieck.umontpellier.fr directly; per the skill transcribe-grothendieck.
% The batch file carries no cover sheet: PDF page k is the archivists' page
% 20+k, and their pencil numbers bottom-left agree throughout. The folder is
% seventy-one pages; this is the second of four batches (the others are
% transcribed in parallel by other passes).
%
% Pass: Opus 5.5 (claude-opus-5-5), 2026-09-26 — first pass, unchecked against the pages by a human.
%
% The folder sits in the inventory group "4-9" « Cristaux »; it has its own
% date in the catalogue, which \dating{} carries verbatim, with no group
% sentence.
%
% Pages skipped: none. Pages summarised: none. All twenty leaves are his own
% manuscript (black ink on pp. 21-22, blue ink from p. 22 on) and are
% transcribed in full, as far as the hand allows.
%
% His pagination and dates. Two runs.
%   21-34  his pp. 2-15 (numbered at the top centre), the continuation of the
%          talk opened on p. 20 of batch 1 (« Nov. Déc. 196? », « 1.
%          Cohomologie de De Rham des var. alg. », his p. 1): end of §1
%          (the comparison theorem, 4) « conséquences métaphysiques »,
%          5) « critique du point de vue : cohomologie de De Rham », items
%          a) struck, then b), c) after a relettering), §2 « Tour d'horizon
%          sur théories cohomologiques l-adiques » (p. 26), §3 « Revue des
%          essais de définition d'une théorie cohomologique p-adique en car.
%          p > 0. Retour sur De Rham via Washnitzer-Monsky » (p. 32), with
%          1) flat cohomology and 2) « L'approche de Washnitzer-Monsky »
%          (p. 34). His p. 15 (our p. 34) breaks off after the boxed
%          H_DR(X/S).
%   35-40  unpaginated; p. 35 is dated at the top right « ?5.11.66 » (the
%          first digit is written over, 1 and 2 superposed; read 15 and
%          flagged). A fresh start of the same material (« Nous avons montré
%          que la cohomologie ... »), then « L'idée de W-M. » (pp. 36-37,
%          remarks a)-c)), then a second restart headed « L'idée de W.M. »
%          (p. 38) which computes H^1 of the formal affine line and
%          introduces the overconvergent ring W^+{t}; p. 40 breaks off
%          mid-sentence (« (les idéaux sont-ils ») into batch 3.
%   The date 15.11.66 bears on batch 1's reading of p. 20 (« Nov. Déc.
%   196? »): these notes look like November 1966, not 1967, though the
%   catalogue title says 1967. Recorded in a \note on p. 35; \dating{} is
%   the catalogue's.
%
% What the batch is about. Notes for a talk (« laïus ») on De Rham
% cohomology and the search for a p-adic cohomology. End of §1: the
% comparison theorem H^*(X^an, C) = H^*_DR(X) for X locally of finite type
% over C, the affine corollary, proof indications (GAGA for proper X,
% resolution of singularities in general); its « metaphysical »
% consequence (transcendental cohomology has a purely algebraic
% description). Critique of De Rham: excellent in char. 0; in char. p
% infinite-dimensional for non-proper X, wrong Betti numbers (dim H^1_DR >=
% 2 dim Pic), coefficients of char. p so no Lefschetz traces; tied to
% smooth varieties (nodal curves); need for Leray spectral sequences,
% f_!, f^! à la Hartshorne, more general coefficients (differential
% operators). §2: l-adic cohomology, comparison H^i(X, Z_l) = H^i(X^an, Z)
% (x) Z_l, independence of l known only via char. 0; in char. p the
% missing p-adic theory, Serre's argument that no Z_p/Q_p theory exists,
% W(k) as the natural coefficient ring. §3: flat cohomology with radicial
% coefficients (Artin: vanishing above n+1), then Washnitzer-Monsky:
% lift X_0 to X smooth over W(k), H_DR(X/S), invariance of H_DR (x) K;
% remarks on the formal-scheme version, local invariants, the generic
% fibre. Pp. 38-40: for the formal affine line H^1 is of infinite rank
% over W (integrating sum a_n t^n needs v_p(a_n) - v_p(n+1) -> +infty),
% hence the overconvergence condition |a_n| = O(1) rho^n, rho > 1, and
% the rings W^+{t_1,...,t_n} and W-M-algebras.
%
% Notation fixed once for the batch (continuing batch 1): \mathbb{H} for his
% double-stroked H (hypercohomology); \mathfrak{X} for his script X (formal
% scheme, pp. 38-39); \Omega^{\cdot} for his Omega with a dot, which he
% often underlines (the underline is not reproduced); W^{\dagger} for the
% superscript cross on p. 40 (Monsky-Washnitzer's dagger); X^{\mathrm{an}}
% for his X^an. His underlined words in running text are set in \emph.
%
% Hardest pages and readings to check first: the crowded, cancelled item
% « a) » and the two struck marginal notes of p. 22; the diagonal marginal
% note of p. 31; the interlinear insertions of pp. 25, 29, 30, 32; the
% phrases « Je prouve dans papier bleu » (p. 21) and « ce pb ... » (p.
% 30); the first digit of the date on p. 35; « K corps des fractions de K »
% on p. 40. Apparatus: 114 \ill{} and 51 \uncertain{}.

\documentclass[11pt,a4paper]{article}
% Le préambule commun à toutes les transcriptions du fonds Grothendieck.
%
% Il tient en une idée : **l'incertitude se compose, elle ne se résout pas.**
% Une transcription qui lisse un mot douteux a détruit la seule chose pour
% laquelle on la faisait — un lecteur doit pouvoir savoir, à chaque endroit,
% ce qui était sur le feuillet et ce qui a été ajouté. Les six macros qui
% suivent sont donc l'appareil critique, et non de la décoration.
%
% Les mêmes commandes sont reconnues par `scripts/render.mjs`, qui produit la
% vue de lecture du site. Une macro ajoutée ici doit l'être aussi là-bas,
% faute de quoi le rendu échouera bruyamment — ce qui est voulu : un rendu
% silencieusement faux est pire qu'un rendu absent.

\ProvidesPackage{grothendieck}[2026/08/08 Transcriptions du fonds Grothendieck]

\usepackage{amsmath,amssymb,amsthm}
\usepackage{mathrsfs}
\usepackage{stmaryrd}
% Les diagrammes commutatifs sont partout dans ce fonds : c'est la langue
% même de la géométrie algébrique de Grothendieck, et une transcription qui
% les rendrait en prose ne transcrirait rien.
\usepackage{tikz-cd}
% Français par défaut — c'est la langue des feuillets — mais l'anglais est
% chargé aussi : la traduction et le résumé sont des documents anglais, et la
% typographie française (espaces avant les deux-points) y serait une faute.
% Ces documents basculent par \selectlanguage{english} en tête de corps.
\usepackage[english,french]{babel}
\usepackage{fontspec}
\usepackage{geometry}
\geometry{a4paper,margin=25mm,marginparwidth=28mm}
\usepackage{marginnote}
\usepackage{xcolor}
% ulem, not soul. `soul`'s \st and \ul are built on a character-by-character
% scan that XeLaTeX's font machinery breaks: the document fails with an
% undefined control sequence rather than degrading. ulem does the same two jobs
% and is Unicode-engine safe. `normalem` stops it hijacking \emph, which the
% transcriptions use for Grothendieck's own emphasis.
\usepackage[normalem]{ulem}
\usepackage{hyperref}
\hypersetup{colorlinks=true,linkcolor=black,urlcolor=black,pdfborder={0 0 0}}

% --- Métadonnées du lot -----------------------------------------------------
% Reprises telles quelles par le script de rendu, qui les affiche en tête de
% la vue de lecture. Elles fixent ce que le fichier prétend couvrir : sans
% elles, une transcription flotte sans point d'attache dans un fonds de seize
% mille feuillets.

\newcommand{\folder}[1]{\gdef\grothFolder{#1}}
\newcommand{\batch}[1]{\gdef\grothBatch{#1}}
\newcommand{\leaves}[2]{\gdef\grothFirst{#1}\gdef\grothLast{#2}}
\newcommand{\foldertitle}[1]{\gdef\grothFoldertitle{#1}}
\gdef\grothFolder{?}\gdef\grothBatch{?}\gdef\grothFirst{?}\gdef\grothLast{?}\gdef\grothFoldertitle{}

% --- Le résumé --------------------------------------------------------------
%
% Il ouvre la lecture modernisée, sous le titre « Résumé », et il est composé
% en petit. Ce n'est pas un ornement : c'est le seul passage du document qui
% ne soit pas des mathématiques, et un lecteur doit pouvoir le reconnaître
% avant de l'avoir lu — puis le sauter s'il connaît déjà le sujet, ou s'y
% arrêter s'il ne le connaît pas. Le filet qui le suit marque la reprise du
% corps.

\newenvironment{resume}
  {\small\setlength{\parskip}{0.45em}}
  {\par\medskip\hrule\medskip}

% --- Mots-clés ---------------------------------------------------------------
%
% En anglais, en clôture du résumé de la lecture modernisée : le vocabulaire
% moderne sous lequel chercher ce que les feuillets construisent. Le manifeste
% du site (`npm run manifest`) les extrait pour étiqueter la cote entière —
% cette ligne est la source unique des tags, il n'y a pas de fichier à côté.
\newcommand{\keywords}[1]{\par\smallskip\noindent
  {\footnotesize\textsc{Keywords} --- \textit{#1}}\par}

% --- L'appareil critique ----------------------------------------------------

% Le numéro de feuillet, en marge. C'est l'ancre qui permet de régler un
% désaccord sur une formule en regardant *un* feuillet plutôt qu'une chemise
% de six cents. Le site s'en sert aussi pour tourner les pages du fac-similé
% quand on fait défiler la transcription.
\newcommand{\leaf}[1]{%
  \marginnote{\footnotesize\color{gray!70}\textsf{#1}}%
  \label{leaf:#1}\ignorespaces}

% Illisible : on ne devine pas. Un mot inventé qui se lit comme les autres est
% le pire résultat possible.
\newcommand{\ill}{\textcolor{red!60!black}{[\,\ldots\,]}}

% Lecture douteuse : le mot est donné, et signalé comme incertain.
\newcommand{\uncertain}[1]{\uline{#1}}

% Ajout de l'éditeur — une lettre manquante, un mot sous-entendu.
\newcommand{\add}[1]{\textcolor{blue!50!black}{[#1]}}

% Biffé par Grothendieck lui-même. Ce qu'il a rayé fait partie du document :
% une rature dit souvent où la pensée a changé de direction.
\newcommand{\struck}[1]{\sout{#1}}

% Note du transcripteur, sur l'état matériel du feuillet.
\newcommand{\note}[1]{\par\smallskip{\small\color{gray!60!black}\textit{[#1]}}\par\smallskip}

% Note marginale de Grothendieck, distincte du corps du texte.
\newcommand{\marginal}[1]{\par\smallskip\noindent\hspace{1em}%
  {\small\color{gray!60!black}#1}\par\smallskip}

% --- Titre ------------------------------------------------------------------

\AtBeginDocument{%
  \begin{center}
    {\large\bfseries Fonds Alexandre Grothendieck --- cote n\textsuperscript{o}~\grothFolder}\\[2pt]
    {\itshape\grothFoldertitle}\\[4pt]
    {\small feuillets \grothFirst--\grothLast\ (lot \grothBatch)}\\[2pt]
    {\footnotesize\color{gray!60!black}Transcription établie d'après le fac-similé numérisé
     par l'Université de Montpellier.\\
     Les lectures incertaines sont soulignées, les illisibles notées [\,\ldots\,],
     les ajouts entre crochets.}
  \end{center}
  \vspace{1em}}

\endinput


\folder{7}
\batch{2}
\pages{21}{40}
\dating{1966-1967}
\watermark{Édition de démonstration}
\foldertitle{Notes laïus novembre-décembre 1967 (Cristaux) : notes manuscrites (1966-1967)}

\begin{document}

\page{21}\note{p.\ 2 de l'auteur ; suite de l'exposé ouvert p.\ 20}

Je \uncertain{prouve} dans papier bleu

\emph{Th.} \emph{C'est un isom.} ($X$ \struck{\ill{}} schéma loc.\ de type
fini \struck{\ill{}} sur $\mathbb{C}$)

\add{Si $X$ affine, \uncertain{c'est équivalent au} :}\note{ajouté au-dessus
de la ligne}

\emph{Corollaire} \struck{\ill{}} Si $X$ affine non sing., alors
$H^{*}(X^{\mathrm{an}}, \mathbb{C})$ peut se calculer avec le complexe de
De Rham « algébrique »

N.B. Le théorème résulte facilement des cas particuliers, par suite
spectrale de Leray d'un recouvrement.

Indications sur démonstration. Si $X$ est \emph{propre}, cela résulte des
suites spectrales de type $(*)$, algébrique et analytique, et de Gaga,
\ill{} \ill{} \ill{} \ill{} \ill{} \ill{} termes initiaux. [C'est chose
essentiellement connue depuis Gaga de Serre]

Le cas général \ill{} \uncertain{ramène} au cas affine, \ill{} \ill{}
\struck{\ill{}} moyen plus puissant que Gaga, la résolution des
singularités. Donner \struck{\ill{}} une idée des principes de la
démonstration \add{du th.\ de comparaison.}\note{ajouté au-dessous de la
ligne et relié par un trait à ce qui suit}

4) \emph{Conséquences métaphysiques} : La cohomologie complexe (définie par
voie transcendante), interprétée comme « cohomologie de De Rham »

\page{22}\note{p.\ 3 de l'auteur}

admet une description purement algébrique, valable \add{pour des schémas
lisses} sur un corps \struck{de car.} quelconque. Elle est \emph{filtrée},
le gradué associé étant les termes $E_\infty$ associés à une suite
spectrale $(*)$. Voir pour des commentaires mon papier bleu\ldots

5) \emph{Critique du point de vue : Cohomologie de De Rham.}

\uncertain{a)} \struck{\ill{} théorie \ill{} \ill{} \ill{} théorie des
\ill{} \ill{} classes de cohomologie associées à des cycles (cf.\ p.\,ex.\
\ill{} coh.\ de Hodge, dans mon \uncertain{ancien} exp.\ Bourbaki) \ill{}
singularités}\note{toute la ligne a) est barrée de traits obliques ; au-dessus,
une insertion à peine lisible, « \uncertain{Le cadre} \ill{}
(\uncertain{non géom.}) », elle aussi soulignée et barrée}

\uncertain{b)}\note{la lettre est recouverte d'une tache d'encre : sans doute
« b) » récrit en « a) » après la suppression de l'item précédent, puisque
l'item suivant, p.\ 23, est « b) » récrit sur « c) »} En car.\ $0$,
\add{pour coh.\ DR.\ des variétés complètes \uncertain{lisses}.}\note{dans
la marge gauche, entouré}
donne des invariants excellents, donnant en particulier les « bons »
nombres de Betti.\marginal{\struck{Mais on sait qu'elle n'est pas à la
hauteur des formalismes généraux en \uncertain{géom.\ alg.} auquel on est
maintenant habitué. Un tel formalisme, par rapport aux \ill{} Hartshorne,
devrait avoir}}\note{dans la marge gauche, encadré et barré de deux traits
obliques}

Mais en car.\ $p$, \struck{\ill{}} n'est \uncertain{aucunement}
raisonnable, car $1^\circ$) si $X$ n'est pas propre, on trouve en général
des modules de dim.\ infinie

\uncertain{voir}
\[
  H^{0}_{\mathrm{DR}}(X) \simeq H^{0}(X, \mathcal{O}_X)^{p} ;
\]
$2^\circ$) Même si $X$ est propre, on ne trouve pas en général les « bons »
nombres de Betti, par exemple on a
\[
  \dim H^{1}_{\mathrm{DR}}(X) \geq 2 \dim \mathrm{Pic}_X ,
\]
et il y a des exemples où on a $>$, (dû « moralement » à des phénomènes de
$p$-torsion\ldots).

Enfin $3^\circ$) Même dans des cas (tels celui des VA) où les
$\dim H^{i}_{\mathrm{DR}}(X)$ \struck{sont $b_i(X)$} sont les « bons »
$b_i(X)$, un défaut majeur de ces \uncertain{espaces} (déjà signalé par
Serre) c'est que le corps de

\page{23}\note{p.\ 4 de l'auteur}

définition est de car $p > 0$, et donc ils ne se prêtent pas à la
détermination de nbs \emph{exacts} de pts fixes par formules de Lefschetz
(mais seulement les nbs mod $p$), donc non plus à une formulation
satisfaisante des conjectures de Weil.

b)\note{la lettre est récrite, sur un « c) » semble-t-il} Le point de vue De
Rham classique, \add{pour exprimer la coh.\ « complexe »,} est trop
exclusivement lié aux variétés \emph{non} singulières. Rien n'empêche
évidemment de considérer aussi le complexe $\Omega^{\cdot}_X$ et
$\mathbb{H}^{*}(X, \Omega^{\cdot}_X)$ si $X$ est quelconque. Mais si $X$
est une courbe, affine \struck{ou projective} ou complète, sur
$\mathbb{C}$, \add{p.\,ex.\ la droite affine ou projective avec un pt
double ordinaire} \struck{intègre \ill{} [\ill{} \ill{} \ill{} \ill{}
\ill{}]}, on n'obtient déjà pas pour $\dim H^{1}_{\mathrm{DR}}(X)$ le
premier nb de Betti comme il faudrait. On peut décider qu'on ne
s'intéresse qu'aux variétés algébriques non singulières --- mais
l'expérience apprend que \struck{\ill{}} pour étudier celles-ci, on ne peut
éviter l'introduction de variétés singulières, p.\,ex.\ à titre de fibres
intervenant dans des \struck{fibrations} morphismes

\page{24}\note{p.\ 5 de l'auteur}

de schémas non lisses : \struck{\ill{}}
\[
  f : X \to Y .
\]
Ainsi, connaissant déjà $H^{*}(X^{\mathrm{an}}, \mathbb{C})$ comme
$H^{*}_{\mathrm{DR}}(X)$, on aimerait une \struck{interprétation}
\add{construction} purement algébrique de la suite spectrale de Leray de
$f^{\mathrm{an}} : X^{\mathrm{an}} \to Y^{\mathrm{an}}$ :
\[
  H^{*}(X^{\mathrm{an}}, \mathbb{C}) \Longleftarrow
  E_2^{pq} = H^{p}\bigl(Y^{\mathrm{an}},
  R^{q} f^{\mathrm{an}}_{*}(\mathbb{C}_{X^{\mathrm{an}}})\bigr) .
\]
En fait, si $f$ est lisse, donc à fibres non singulières, une telle
description \uncertain{existe} par des \ill{} tout clair en termes de ce
qui a été \ill{} jusqu'à maintenant.
\struck{[les $R^{q} f_{*}(\Omega^{\cdot}_{X/\mathbb{C}})$]} \add{Nous y
reviendrons en} \ill{} \ill{} : l'\uncertain{heure}. Un formalisme
\struck{\ill{} des faisceaux \ill{} sur $Y$, par des}
\add{cohomologique}\note{entouré, dans la marge, relié par un trait}
\add{vraiment} maniable devrait impliquer, tout comme dans le cadre
classique, l'écriture de suites spectrales de Leray --- ce qui n'est
possible que s'il englobe des « faisceaux » de coefficients de nature plus
générale que les complexes de De Rham $\Omega^{\cdot}_X$, --- \add{englobant,
sans doute, également des complexes d'opérateurs différentiels de nature
plus générale que De Rham.}

c) Du point de vue dualité et notions \uncertain{similaires},

\page{25}\note{p.\ 6 de l'auteur}

notons que pour des $X$ lisses et propres, on a bien le formalisme de
dualité habituel (autodualité de $H_{\mathrm{DR}}(X)$, \add{d'où} th.\ de
Gysin, classe de coh.\ associée à un cycle alg., compatibilité avec images
directes, inverses, intersections)\ldots{} la chose se faisant de façon tout
analogue à ce que, dans le \uncertain{cadre} de la « cohomologie de Hodge »
j'avais fait au Sém.\ Bourbaki. Mais le « lecteur informé » au courant des
histoires $f_{!}$, $f^{!}$ notera à quel point une \add{simple} dualité de
ce type est \uncertain{étriquée} et \uncertain{pauvre}. On aimerait un
formalisme $f_{!}$, $f^{!}$ du type de celui développé en long et en large
dans le Sém.\ Hartshorne (Springer), \struck{\ill{}} pour des morphismes
\add{disons \struck{propres,} séparés de type fini,} quelconques, de
schémas. Un tel formalisme \add{(impliquant des coeff.\ plus généraux que
les complexes de DR.)} devrait \struck{\ill{}} automatiquement donner
également une description de l'\emph{homologie} singulière (comme il est
suggéré par ailleurs dans mes papiers bleus\ldots). Il y aurait lieu en
particulier d'étendre ce formalisme, du cas de complexes à différentielles
linéaires, au cas des complexes

\page{26}\note{p.\ 7 de l'auteur}

à op.\ différentiels.\note{au-dessus, relié par un trait :
« différentielles des », de sorte qu'on lit peut-être « complexes à
différentielles des op.\ différentiels »}

Dans b) et c) sont posées, essentiellement, la question de la définition
de la bonne \emph{catégorie de coefficients}, pour englober à la fois le
formalisme des complexes \struck{\ill{}} « linéaires » cohérents, et des
complexes du type De Rham. \struck{Les} exposés que je fais doivent donner
\struck{quelques} des indications plus précises sur \struck{la} une voie
\struck{dans} \add{\uncertain{suivie}} dans laquelle il conviendrait
\struck{de} probablement de chercher des catégories de coefficients.

\section*{§\,2 Tour d'horizon sur \struck{la} \struck{\ill{}} théories
cohomologiques \add{$\ell$-adiques} pour les variétés algébriques
abstraites.}

\struck{\ill{}}\note{un mot biffé, noirci, après le titre}

On peut évidemment se demander : pourquoi \uncertain{à tout prix}
\struck{\ill{}} des progrès \uncertain{sont} dans nos programmes
d'algébrisation \add{et de généralisation} à outrance de la cohomologie de
De Rham. On dispose déjà des invariants cohomologiques $\ell$-adiques, qui
se prêtent à des applications \struck{géométriques}

\page{27}\note{p.\ 8 de l'auteur}

arithmétiques comme la formulation cohomologique des conjectures de Weil.
Faisons un petit tour d'horizon \ill{} \ill{} on en est \uncertain{aussi}
à ce sujet.

Si $X$ est une variété alg.\ sur un corps alg.\ clos $k$,
\struck{\uncertain{propre} sur $k$ \ill{} \ill{} \ill{} \ill{} les idées.}
\add{\uncertain{on associe} pour \struck{\ill{}} \ill{} $\ell$ nb premier}
\emph{$\ell \neq \mathrm{car}\, k$} des modules sur $\mathbb{Z}_\ell$
\[
  H^{i}(X, \mathbb{Z}_\ell) = \varprojlim_{\nu}
  H^{i}(X_{\mathrm{ét}}, \mathbb{Z}/\ell^{\nu}\mathbb{Z}) ,
\]
dont les \uncertain{comportements} pour $\ell$ \struck{variable} fixé, $X$
variable, \ill{} à peu près celui de la cohomologie \add{entière}
\uncertain{traditionnelle} dans le cadre transcendant. Ainsi, si $k$ est le
corps des complexes, le th.\ de comparaison nous fournit des isom.\ can.,
fonctoriels en $X$
\[
  \boxed{\;H^{i}(X, \mathbb{Z}_\ell) \simeq
  H^{i}(X^{\mathrm{an}}, \mathbb{Z}) \otimes_{\mathbb{Z}} \mathbb{Z}_\ell\;}
\]
Ceci montre que \emph{la connaissance de la cohomologie entière
transcendante de $X^{\mathrm{an}}$} (du moins du point de vue additif) est

\page{28}\note{p.\ 9 de l'auteur}

\struck{entièrement déterminée par} \add{équivalente à} : \emph{la
connaissance de tous les $H^{i}(X, \mathbb{Z}_\ell)$}.\note{« tous » est
entouré} De façon plus précise, pour un $\ell$ donné, la connaissance de
$H^{i}(X, \mathbb{Z}_\ell)$ équivaut à celle du rang \add{$b_i(X)$} de
$H^{i}(X^{\mathrm{an}}, \mathbb{Z})$, \add{égal à celui de
$H^{i}(X, \mathbb{Z}_\ell)$ sur $\mathbb{Z}_\ell$)} et de la composante
$\ell$-primaire du sous-groupe de torsion de
$H^{i}(X^{\mathrm{an}}, \mathbb{Z})$. On obtient, en passant, par voie
transcendante, le fait que si le corps de base $k$ est de car.\ $0$, alors
les \add{$H^{i}(X, \mathbb{Z}_\ell)$ sont de t.f., leurs rangs}
$b_i(X, \mathbb{Z}_\ell)$ \struck{sont} pour $\ell$ variable est
indépendant de $\ell$, et que pour $X$ fixé, sauf pour un nb fini de
$\ell$, \struck{le} $H^{i}(X, \mathbb{Z}_\ell)$ \ill{} \ill{} torsion. Et
par la même voie, on trouve \add{p.\,ex.} que si on a un endom.\ $f$ de
$X$, alors pour $i$ fixé, \struck{les coeff.} le polynôme caractéristique
de l'endom.\ de $H^{i}(X, \mathbb{Z}_\ell)$

\page{29}\note{p.\ 10 de l'auteur}

induit par $f$ est \emph{à coefficients entiers indépendants de $\ell$}.

Situation beaucoup moins bonne en car.\ $p > 0$, où on ne dispose pas de
« cohomologie entière » transcendante. Le fait que les
$H^{i}(X, \mathbb{Z}_\ell)$ sont de t.f.\ (résultat de ce que les
$H^{i}(X, \mathbb{Z}/\ell\mathbb{Z})$ le sont) n'est prouvé que si $X$
est propre ou $\dim X \leq 2$ --- pour l'avoir en général, il semble qu'il
faille la résolution des singularités. D'autre part, \struck{on} si $X$ est
projectif et lisse, on ignore si p.\,ex.\ le rang de
$H^{2}(X, \mathbb{Z}_\ell)$ est indépendant de $\ell$, et à fortiori si le
résultat plus fort pour les coefficients des polynômes caractéristiques
\add{\uncertain{sont}} valable. Cela serait \add{d'ailleurs},
actuellement, \struck{le premier} \ill{} plus important \ill{} en
direction d'une solution des conjectures de Weil. \struck{\ill{}} On peut
commencer par tensoriser par $\mathbb{Q}_\ell$,

\page{30}\note{p.\ 11 de l'auteur}

d'où des coefficients dans les $\mathbb{Q}_\ell$, et un premier problème
est s'ils sont dans $\mathbb{Q}$. On sait (si $X$ est projective non
singulière) \uncertain{ramener} ce pb à un pb de cycles algébriques du type
Lefschetz.

Mais --- c'est là le « point » que je désire faire --- même si ce pb
« \ill{} » \add{i.e.\ « mod torsion »} est résolu par l'affirmative, on ne
pourra aucunement conclure que les coeff.\ sont des entiers, mais
seulement \add{qu'ils sont dans $\mathbb{Z}[1/p]$ i.e.} \struck{des
entiers $\ell$-adiques} de la forme $n/p^{s}$, avec $n \in \mathbb{Z}$.
Ceci provient de la restriction
\[
  \ell \neq \mathrm{car}\, k ,
\]
et montre qu'il \emph{manque une théorie cohomologique $p$-adique pour les
variétés algébriques en car.\ $p$}. Faute d'avoir une théorie de ce type,
on ne saura

\page{31}\note{p.\ 12 de l'auteur}

ni plus comprendre les \emph{phénomènes de $p$-torsion cohomologique} en
var.\ alg.\ en car.\ $p$. Des \struck{\ill{}} motivations \add{pour
\uncertain{vouloir} \ill{} \uncertain{complète}} ne manquent pas, cf
p.\,ex.\ l'exposé de Tate sur Brauer, \struck{\ill{}} \add{on le
\uncertain{établissait}} \struck{par} heuristique d'une formule pour
?\marginal{\ill{} dans \ill{} Brauer d'un \ill{} \uncertain{sur un} corps
fini \ill{} \ill{} \ill{} \ill{} \ill{} \ill{} « \ill{} $p$-adiques »
\ill{}\ldots}\note{dans la marge gauche, écrit en oblique, relié au texte ;
presque entièrement illisible} \struck{En analogie avec le cas
$\ell$-adique, on} \add{\ill{} \ill{}} espère trouver un foncteur
\[
  X \longmapsto H^{\cdot}(X)
\]
de la catégorie des VA projectives lisses sur $k$, dans celle des modules
de t.f.\ sur $\mathbb{Z}_p$, ou sur $\mathbb{Q}_p$, ayant les propriétés
habituelles \struck{cf} (Künneth, dualité, « bons » nbs de Betti\ldots). Un
argument de Serre montre que déjà pour le $H^{1}$, et en se bornant aux
$X$ \struck{de $\dim \leq 2$}, \uncertain{qui sont des} VA de
$\dim \leq 2$, un tel foncteur n'existe pas. Il semble que l'anneau de
coefficients naturel d'une coh.\ « $p$-adique », pour les var.\ sur un
corps $k$ de car.\ $p > 0$, \struck{\ill{}} \ill{} l'anneau

\page{32}\note{p.\ 13 de l'auteur}

$W(k)$ des vecteurs de Witt sur $k$, anneau de val.\ discrète dont le corps
des fractions $W(k)[\frac{1}{p}] \supset \mathbb{Q}_p$ est un corps de
car.\ \emph{nulle}.

Ceci devient \uncertain{clair} \ill{} \ill{} \ill{} \ill{} un corps de base
\add{parfait} quelconque, --- par voie \uncertain{alg.} \ill{}\ldots

\section*{§\,3 Revue des essais de définition d'une théorie
\add{cohomologique} $p$-adique en car.\ $p > 0$. Retour sur De Rham via
Washnitzer-Monsky. \add{Lien avec les stratifications de Gauss-Manin.}}

1) \emph{Utilisation des \uncertain{$p$-groupes} éventuellement radiciels},
\struck{pour} comme coefficients pour la cohomologie (avec la topologie
plate).

Heuristiquement, cela semble magnifique \add{\ill{} \ill{} une variété}
en dim.\ $1$, et \struck{\ill{} \ill{} \ill{}}
\add{considérablement \uncertain{moins}, de nature géométrique,}
\struck{\ill{} \ill{} \ill{}} en termes des théorèmes de dualité en
cohomologie plate, en termes de classes \ill{} \ill{} géométriques
(courbes sur un

\page{33}\note{p.\ 14 de l'auteur}

corps, éventuellement fini). On espère que Raynaud tirera au clair cette
théorie, et ses analogues locaux --- ce qui exigera d'ailleurs de surmonter
des obstacles techniques considérables.

Mais en dim $n \geq 2$, c'est sans espoir, car Artin a remarqué que tous
les \struck{groupes de} coh.\ de $X$, à coeff.\ dans des groupes radiciels,
s'annulent en dim $> n+1$. Donc \emph{on \ill{} les classes \ill{}}, \emph{on
renonce à construire les « bons » invariants cohomologiques par cette voie
\ill{} \ill{}} $\dim X + 1$ (et non $2 \dim X$, qui est la bonne). Il
faudra donc (si $X$ projective lisse) \add{\ill{} \ill{} mais} compléter
le tableau par heuristique par dualité de Poincaré\ldots{} On voit donc
qu'il faudrait

\page{34}\note{p.\ 15 de l'auteur}

\emph{une définition de la coh.\ $p$-adique indépendante de la coh.\
plate}, \emph{quitte ensuite} à étudier ses relations (sans doute fort
intéressantes) avec cette dernière.

2) \emph{L'approche de Washnitzer-Monsky.}

C'est dans \struck{cette} l'approche \uncertain{W-M} qu'on voit la
cohomologie de De Rham \struck{qui} montrer à nouveau le bout de
l'oreille\,! Nous allons expliquer \add{pour le moment} comment elle
s'introduit dans l'approche de W-M., en \struck{\ill{}} simplifiant leur
construction à l'extrême, (au risque de la défigurer\ldots\,!). On part de
$X_0$ \struck{projectif} lisse sur $s$, on \uncertain{suppose}
\add{d'abord} qu'on puisse relever en $X$ \struck{projectif} propre et
lisse sur $S = \operatorname{Spec} W(k)$, \struck{Soit \ill{} les}
\struck{point générique} Considérons le complexe $\Omega^{\cdot}_{X/S}$, et
\[
  \boxed{\;H^{\cdot}_{\mathrm{DR}}(X/S) =
  \mathbb{H}^{*}(X, \Omega^{\cdot}_{X/S})\;}
\]
C'est un module de t.f.\ sur $W(k)$.

\begin{tikzcd}
X_0 \arrow[r, hook] \arrow[d, no head] & X \arrow[d, no head] \\
s = \operatorname{Spec} k \arrow[r, hook] & S = \operatorname{Spec} W
\end{tikzcd}

\note{ce carré est dessiné dans la marge gauche, en face de « On part de
$X_0$ \ldots{} Considérons le complexe » ; les traits verticaux sont sans
flèche, les inclusions notées $\subset$}

\page{35}

\marginal{\uncertain{15}.11.66}\note{en haut à droite ; le premier chiffre
est récrit, un 1 et un 2 superposés. Cette date suggère que la date
« Nov.\ Déc.\ 196? » de la p.\ 20 (batch 1) est bien 1966, bien que le
titre du catalogue porte 1967}

Nous avons montré que la cohomologie \struck{de DR.} \add{complexe
ordinaire} d'une variété algébrique non singulière sur $\mathbb{C}$ se
décrit par voie purement algébrique comme « cohomologie de De Rham »
\[
  H_{\mathrm{DR}}(X) = \mathbb{H}^{*}(X ; \Omega^{\cdot}_X)
\]
Celle-ci a un sens pour n'importe quel schéma sur un corps quelconque.
Cependant cette version de la cohomologie a de nombreux désavantages : liée
aux formes diff.\ donc aux variétés non sing., \add{le corps de
\struck{coefficients} base de la théorie coh.\ peut être de car $p > 0$\ldots}
D'autre part, nous avons vu la désirabilité d'avoir une théorie de
cohomologie « $p$-adique » pour les variétés \struck{en} en car.\ $p$, pour
jouer le même rôle que les cohomologies $\ell$-adiques, $\ell \neq p$. On
veut les \struck{\ill{}} nbs de Betti etc\ldots

Une remarque sur la définition de la c.\ de De Rham :
$\Omega^{1}_{X/S}$ est défini pour un préschéma relatif $X/S$, et
$\bigwedge^{\cdot} \Omega^{1}_{X/S} = \Omega^{\cdot}_{X/S}$ est muni d'une
op.\ diff.\ naturelle\ldots{} On peut donc définir
\[
  \left\lbrace
  \begin{array}{l}
    H_{\mathrm{DR}}(X/S) = \mathbb{H}^{*}(X, \Omega^{\cdot}_{X/S}) \\
    \text{et les } Rf_{*}(\Omega^{\cdot}_{X/S})
  \end{array}
  \right.
\]
\add{[coh.\ \uncertain{si} $f$ propre et lisse, \ill{}]}\note{en bout de
ligne, à droite de la seconde formule, en petits caractères}
\marginal{\ill{} qui sont \uncertain{comp.}\ \ill{} avec \uncertain{chgt de
base} si $f$ \emph{lisse}\ldots}\note{en bas à gauche, encadré}

\page{36}

\emph{L'idée de W-M.} \struck{serait} \add{semble pouvoir s'exprimer} plus
ou moins \add{en disant} que \emph{ce module est un invariant qui, à isom.\
canonique près, ne dépend pas du choix du relèvement $X$ de $X_0$, et qui
serait donc un invariant intrinsèquement attaché à $X_0$, qui serait alors
par définition la cohomologie \struck{de} « $p$-adique » cherchée}.
\marginal{« \uncertain{Some properties of} formal schemes » 63/64\ldots}\note{dans
la marge gauche, écrit en oblique}

\emph{Remarques} a) Cette construction \uncertain{marcherait} aussi,
\uncertain{mutatis mutandis}, si on ne peut relever $X_0$ qu'en un
\emph{schéma formel} lisse sur $S$.

b) D'après le texte publié de W-M, \struck{il ne semble pas qu'il} on
pourrait \struck{\ill{}} être amené à penser que seul
$H^{\cdot}_{\mathrm{DR}}(X/S) \otimes_W K$ \add{$K$ corps des fractions de
$W$} devrait avoir la propriété d'invariance en question. De plus, la
construction et la propriété invariance n'y est traitée que pour $X_0$
\emph{affine}

\page{37}

(et non propre), en utilisant un relèvement de $X_0$, non comme schéma ni
comme schéma formel, mais en quelque chose qui ressemble à une
\emph{variété rigide-analytique} à la Tate. De ce fait, ils n'ont défini,
\add{dans [~]} que des invariants \emph{locaux}. Par communication orale de
W., il serait aisé de définir les invariants globaux (sur $W$, et non
seulement sur $K$), en travaillant \add{avec le complexe de De Rham} sur le
\emph{site} des \emph{relèvements de W.M} \struck{\ill{}} des \emph{ouverts
affines}, et sans avoir à justifier les \uncertain{assertions} que j'ai
faites.

c) On sait, grâce au th.\ de comparaison, que par cette construction
(quand $X_0$ \struck{se} admet un $X$ lisse sur $S$) on trouve des
$H^{\cdot}(X_0)$ qui donnent les bons nbs de Betti --- ce sont
\struck{ceux} ceux de la fibre générique $X_t$\,!

\page{38}

\section*{L'idée de W.M.}

Donnons quelques détails sur ce que font \emph{réellement} W.M. Faute de
pouvoir relever en général un $X_0$ lisse quelconque, on travaille
\add{d'abord} avec des $X_0$ \emph{affines}. On sait (SGA 1960/61, III)
qu'un tel $X_0$ se relève en un $\mathfrak{X}$ \emph{formel} \add{lisse}
sur $S = \operatorname{Spec} W(k)$, \add{unique à isom.\ (\emph{non
unique}) près.} La première idée serait de prendre
$\Omega^{\cdot}_{\mathfrak{X}/S}$ défini par passage à la limite en $n$ à
partir des $\Omega^{\cdot}_{X_n/W_n}$, et de poser
\[
  H(X_0) = \mathbb{H}^{*}(\mathfrak{X}, \Omega^{\cdot}_{\mathfrak{X}/S})
  = H^{*}\bigl(\Gamma(\mathfrak{X}, \Omega^{\cdot}_{\mathfrak{X}/S})\bigr)
\]
Mais montrons que cela ne donne pas du tout ce qu'il faut, \struck{\ill{}}
si $X_0 = \operatorname{Spec} k[t]$ est la droite affine, donc
\[
  \mathfrak{X} = \operatorname{Spf} W\{t\} ,
  \qquad W\{t\} = \varprojlim W_n[t]
\]
\note{sous $W\{t\}$, relié par une accolade : « séries formelles
restreintes, i.e.\ à coeff.\ $\to 0$ »}
On voudrait trouver en effet que
\[
  \mathbb{H}^{1}(\mathfrak{X}, \Omega^{\cdot}_{\mathfrak{X}/S}) = 0
\]
\note{l'exposant $1$ est récrit sur un autre chiffre}

\page{39}

ou tout au moins que le premier mb.\ soit de torsion. En fait, il est de
\emph{rang infini} sur $W$\,!
\[
  \begin{aligned}
  \mathbb{H}^{1}(\mathfrak{X}, \Omega^{\cdot}_{\mathfrak{X}/S})
  &= \Omega^{1}(\mathfrak{X}/S) / d\bigl(\mathcal{O}(\mathfrak{X}/S)\bigr) \\
  &\simeq W\{t\} / \mathrm{Im}\,(f \mapsto f')
  \end{aligned}
\]
\note{sous $\mathrm{Im}(f \mapsto f')$, relié par une accolade : « séries
formelles restreintes \emph{dérivées} »}

Or pour que $\sum a_n t^{n}$ soit une dérivée de $\sum b_n t^{n}$, i.f.\ et
s.\ que
\[
  \sum \frac{1}{n+1}\, a_n t^{n+1}
  \quad \Bigl(\text{ou encore } \sum \frac{1}{n}\, a_{n-1} t^{n}\Bigr)
\]
soit encore une série formelle restreinte, i.e.\ que
\[
  v_p(a_n) - v_p(n+1) \longrightarrow +\infty
\]
et comme $v_p(n)$ (pour $n$ variable) prend des valeurs arbitrairement
grandes, on voit qu'il ne suffit pas que
\[
  v_p(a_n) \longrightarrow +\infty\,!
\]
En notations multiplicatives, si $|\ |_p$ est la valeur absolue
$p$-adique, \struck{on} veut que
\[
  \frac{|a_n|_p}{|n+1|_p} \longrightarrow 0
\]
donc que
\[
  |a_n|_p = \varepsilon_n\, |n+1|_p , \qquad \varepsilon_n \longrightarrow 0
\]

\page{40}

C'est là une condition \add{de divisibilité des coefficients} arithmétique
assez biscornue. W.M. la remplacent par une condition plus forte, et plus
maniable, savoir que
\[
  |a_n| = O(1)\, \rho^{n} \qquad \text{pour un } \rho > 1
\]
convenable (dépendant de la \uncertain{suite} $(a_n)$).

Cela signifie qu'au lieu de regarder l'anneau $W\{t\}$ des séries
convergentes \add{dans $\overline{K}$ ($K$ corps des fractions de
\uncertain{$K$})}\note{sic : on attendrait « de $W$ »} sur le disque
unité, ils regardent l'anneau $W^{\dagger}\{t\}$ des séries convergentes
sur un disque de rayon $\rho > 1$ \struck{(\ill{} \ill{})} un peu plus grand
(dépendant de la fonction envisagée). On définit de même
$W^{\dagger}\{t_1, \ldots, t_n\}$, et les \struck{\ill{}} W-M-algèbres,
\struck{les W.M} qui sont les quotients de tels anneaux par des idéaux
\struck{quelconques} \ill{} quelconques \struck{\ill{}} (les idéaux
sont-ils\note{la phrase se poursuit dans le batch suivant}

\end{document}
